\section{Experimental Reinterpretation}


This section uses the log-referencing model formalized in preceding sections to reinterpret representative existing consciousness experiments.

The aim is not to deny the experimental results.

Rather, it is to re-read what the existing experiments were showing — not through a consciousness-generation model, but through a log-referencing model across time scales.

Much of the confusion in consciousness research arises from placing processing, action, report, and consciousness on a single timeline.

Under that view, whatever appears earlier tends to be interpreted as cause, and whatever appears later as effect.

But in the framework of this paper, the system does not operate on a single clock.

The fast layer predicts, the slow layer sediments, and the plasticity gate introduces delay and viscosity between the two.

The time differences and report-inaccessibility observed in experiments therefore do not necessarily indicate the absence of consciousness.

They indicate the temporal and pathway conditions required for processing to reach a log-referenceable state.

\subsection{Libet-Type Experiments}


In Libet-type experiments, subjects are asked to move a wrist or finger at a time of their choosing and to report the moment at which they "decided to move." Brain activity is recorded simultaneously.

As is well known, a readiness potential rises in advance of the movement, and the subject's conscious report of intention appears after it.

Traditionally this result has often been interpreted as: the brain decides first and consciousness learns of it afterward, or alternatively, free will is an illusion.

But this interpretation presupposes that decision, brain activity, and conscious report are placed on the same timeline.

This paper does not adopt that presupposition.

The readiness potential can be read as the formation of an action candidate in the fast layer — a signal that the future-oriented predictive and motor dynamics have begun falling into a certain action attractor.

The report of "I just decided to move," on the other hand, appears after that fast processing has connected to the log-referencing structure of the slow layer.

Only at that point does processing take the reportable form of "I am about to do this."

The time difference observed in Libet-type experiments is therefore not a delay in consciousness but the transformation time between speed scales.

Writing the fast-layer processing speed as $v_f$ and the slow-layer sedimentation speed as $v_s$:

\begin{equation}
\Delta v = v_f - v_s
\end{equation}

This speed difference appears as internal delay $\tau$:

\begin{equation}
\tau \propto \frac{\Delta v}{\mu}
\end{equation}

Here $\mu$ is the coupling viscosity between fast and slow layers.

What Libet-type experiments measured was therefore not the presence or absence of free will, but the delay required for fast action preparation to be transformed into a reportable log in the slow layer.

In this view, even the expression "consciousness is delayed" is not quite right.

Consciousness does not operate on the same clock as the fast layer to begin with.

Consciousness opens at the point at which fast processing has been transformed into a log-referenceable state.

For this reason, the earlier appearance of the readiness potential does not show that consciousness is powerless.

It shows that consciousness is not action preparation itself, but the passage point at which action preparation sediments into the slow layer and becomes referenceable.

\subsection{Why ``Free Will'' Is an Ill-Posed Question Here}


The reason Libet-type experiments are repeatedly drawn into discussions of free will is the presupposition that "decision" occurs at a single point.

But in the structure of this paper, decision does not localize to a single moment.

In the fast layer, multiple action candidates arise predictively.

The plasticity gate regulates which of them can sediment into the slow layer.

In the slow layer, past logs, meaning gradient, bodily state, and context compose the reportable form.

Compressing this whole into a single "decision point" loses the temporal structure.

It is therefore structurally inappropriate to ask of Libet-type experiments whether "consciousness decided or the brain decided."

The appropriate questions are: which processing preceded in the fast layer; through what delay $\tau$ did that processing connect to the slow layer; which logs were referenced by the meaning gradient $\nabla\Phi$; and at what point did a reportable form obtain.

Reframing the questions this way, Libet-type experiments are not experiments that deny free will — they are experiments that make visible the limits of the single-timeline model.

\subsection{Split-Brain Experiments}


In split-brain experiments, the corpus callosum is severed, restricting information sharing between the left and right hemispheres.

In this situation, appropriate behavioral responses can occur to stimuli presented to the right hemisphere while verbal report of those stimuli may be impossible.

This phenomenon has often been described as evidence for "two consciousnesses" or "two selves."

But in the framework of this paper, this reading is excessive.

What split-brain experiments show is not the doubling of consciousness.

They show that the pathways of processing, sedimentation, and report can become asymmetric.

Even if processing occurs in the right hemisphere, if that processing does not reach the left-hemisphere-dominant verbal reporting pathway, the subject cannot explicitly articulate it.

This does not mean the processing is absent.

The fact that behavioral responses occur shows clearly that local processing is running.

But that processing has not been integrated into the reportable slow log.

What is lost in split-brain experiments is therefore not processing itself, but the referencing pathway.

This point corresponds to Libet-type experiments.

In Libet-type experiments, a temporal delay appears between processing and report.

In split-brain experiments, a pathway disconnection appears between processing and report.

In both cases, the evidence shows that explicit consciousness requires log sedimentation into the slow layer and referenceability.

\subsection{Reporting Bottleneck and Log Accessibility}


What is particularly important in split-brain experiments is not conflating reportability with consciousness.

Reportability is the pathway by which part of the conscious state is externalized.

But inability to report does not mean processing has not occurred.

In the framework of this paper, explicit report requires at least three conditions:

First, processing has occurred in the fast layer.

Second, that processing has sedimented into the slow layer and become log-referenceable.

Third, that log is connected to the verbal or behavioral reporting pathway.

In split-brain patients, this third condition becomes asymmetric.

As a result, certain processing appears as behavior but not as verbal report.

There is no need to read this asymmetry as "two consciousnesses."

More precisely, part of the log-referencing pathway and reporting pathway has been severed.

\subsection{Rationalization as Log Completion}


A well-known phenomenon in split-brain experiments is that the left hemisphere provides retrospective explanations for actions it could not actually access.

This is often called "confabulation" or "rationalization."

In the position of this paper, this rationalization is not deliberate falsification.

It is the natural behavior of the slow log-referencing system attempting to maintain coherence and fill in missing logs.

The slow layer maintains self-continuity using referenceable logs.

But when the processing pathway is severed, actions occur whose generating logs cannot be referenced.

In this situation, the system generates the most coherent explanation available from referenceable logs rather than leaving the gap as a gap.

This rationalization does not show that consciousness is lying.

Rather, it shows that consciousness constructs coherence from referenceable logs.

In this respect, split-brain experiments weaken the view that the ego is a subject with full access.

The ego is not a center that knows everything.

It is a transformation point that assembles self-continuity from referenceable logs.

\subsection{Experimental Correspondence to the Model}


Libet-type experiments and split-brain experiments appear on the surface to address different problems.

The former concerns temporal difference; the latter concerns interhemispheric pathway disconnection.

But in the framework of this paper, both are different cross-sections of the same structure.

Libet-type experiments make visible the following relation:

\begin{equation}
\text{fast preparation} \rightarrow \text{delayed sedimentation} \rightarrow \text{reportable timestamp}
\end{equation}

Split-brain experiments make visible the following relation:

\begin{equation}
\text{localized processing} \rightarrow \text{blocked transmission} \rightarrow \text{asymmetric reportability}
\end{equation}

In both cases, processing itself and conscious report are not identical.

Processing appears as explicit consciousness only after passing through the plasticity gate, delay, sedimentation, and reporting pathway.

Using this structure, the existing experiments are repositioned as follows:

\begin{itemize}
\item Libet-type experiments: exposure of speed-scale difference and internal delay $\tau$.
\item Split-brain experiments: asymmetrization of log-referencing pathway and reporting pathway.
\item Rationalization: slow-layer coherence maintenance in response to inaccessible log gaps.
\item Conscious report: externalization of sedimented logs, not processing itself.
\end{itemize}

Through this repositioning, consciousness is neither the source of causation nor a mere retrospective illusion.

Consciousness is the state that opens when processing has been transformed into a log-referenceable form.

\subsection{What These Experiments Do Not Show}


Finally, what these experiments do not show should be made explicit.

Libet-type experiments did not show that free will does not exist.

They showed that action preparation and conscious report arise on different time scales.

Split-brain experiments did not show that two selves exist inside a person.

They showed that when processing pathways and reporting pathways are severed, log referencing becomes asymmetric.

Rationalization did not show that consciousness always speaks falsely.

It showed that the slow layer fills in gaps from inaccessible logs in a coherent form.

These experiments therefore do not erase consciousness.

Rather, they return consciousness to its correct position: not the origin of action, not an object generated in the brain, but the state in which log referencing across speed scales is open in a reportable form.

\subsection{Summary}


This section reinterpreted Libet-type and split-brain experiments from the log-referencing model of this paper.

Common to both is the non-identity of processing and conscious report.

Processing runs. But it does not appear as explicit consciousness unless it sediments into the slow layer, becomes referenceable as a log, and connects to the reporting pathway.

What the existing experiments were therefore showing was not the powerlessness of consciousness.

It was the temporal and pathway conditions required for consciousness to obtain.
